\documentclass[conference]{IEEEtran}
\IEEEoverridecommandlockouts
% The preceding line is only needed to identify funding in the first footnote. If that is unneeded, please comment it out.
\usepackage{cite}
\usepackage{amsmath,amssymb,amsfonts,amsthm}
\usepackage{algorithmic}
\usepackage{graphicx}
\usepackage{textcomp}
\usepackage{xcolor}
\usepackage{booktabs}
\usepackage{multirow}
\usepackage{url}
\usepackage{subcaption}

\newtheorem{theorem}{Theorem}
\newtheorem{lemma}{Lemma}
\newtheorem{definition}{Definition}

\def\BibTeX{{\rm B\kern-.05em{\sc i\kern-.025em b}\kern-.08em
    T\kern-.1667em\lower.7ex\hbox{E}\kern-.125emX}}

\begin{document}

\title{Hierarchical Adaptive Fuzzy Control (HAFC) for Geospatial Multi-Zone Precision Irrigation with Hard Mass Conservation and Real-Time Atmospheric Forcing\\
\thanks{This research was developed under the Smart Multizone Agricultural Cyber-Physical Systems Initiative.}
}

\author{\IEEEauthorblockN{1\textsuperscript{st} Pranshu Rajan}
\IEEEauthorblockA{\textit{Department of Electronics and Instrumentation Engineering} \\
\textit{National Institute of Technology / Engineering Division}\\
India \\
pranshu.rajan@ieee.org}
\and
\IEEEauthorblockN{2\textsuperscript{nd} Divyansh Patel}
\IEEEauthorblockA{\textit{Department of Electronics and Instrumentation Engineering} \\
\textit{National Institute of Technology / Engineering Division}\\
India \\
divyansh.patel@ieee.org}
}

\maketitle

\begin{abstract}
Water scarcity imposes an existential bottleneck on global agricultural sustainability, with conventional irrigation relying predominantly on static on-off threshold heuristics or unconstrained classical Proportional-Integral-Derivative (PID) controllers. Such approaches systematically fail to accommodate the non-linear hydraulic transport delays, diurnal microclimate variations, and strict reservoir capacity limits characteristic of multi-zone agro-ecosystems. This paper presents the mathematical formulation, system implementation, and empirical validation of a Hierarchical Adaptive Fuzzy Control (HAFC) cyber-physical platform for multi-zone precision agriculture. The architecture unifies five cascaded Mamdani Fuzzy Inference Systems (FIS) with real-time Penman-Monteith (FAO-56) evapotranspiration energetics, real-time spatial geocoding across fifteen agro-climatic sub-basins, and an analytical bounded priority-weighted water-filling allocation algorithm. The closed-loop cascade decouples root-zone soil hydrology (Stage 1), atmospheric evaporative stress (Stage 2), crop phenological water demand (Stage 3), and supervisory safety throttling (Stage 4), before executing multi-zone supply arbitration (Stage 5) under a provable hard mass conservation invariant ($|\text{Residual}| < 10^{-6}\text{ mm}$). Integrated with an interactive WebGIS dashboard, spatial region selection triggers end-to-end inference in under $18.2\text{ ms}$. Long-term agronomic simulations across 96 continuous hours demonstrate that HAFC achieves a 49.72\% reduction in water consumption compared to on-off controllers and a 12.36\% reduction compared to tuned PID feedback, while suppressing root-zone moisture tracking Root Mean Square Error (RMSE) to 0.42\% and completely eliminating transient crop moisture stress days.
\end{abstract}

\begin{IEEEkeywords}
Precision Agriculture, Hierarchical Adaptive Fuzzy Control (HAFC), Cyber-Physical Systems (CPS), FAO-56 Penman-Monteith, Bounded Water-Filling Optimization, Mass Conservation Invariant, Geospatial Telemetry, Soil Hydrology.
\end{IEEEkeywords}

\section{Introduction}
\IEEEPARstart{A}{griculture} accounts for more than 70\% of global freshwater withdrawals, with significant volumetric inefficiencies stemming from archaic hydraulic control paradigms \cite{allen1998crop}. In developing and agrarian economies such as India, escalating climate volatility, erratic monsoon onset, and over-drafted aquifers have created an acute need for autonomous, closed-loop precision irrigation.

Conventional automated irrigation deployments in commercial fields suffer from severe structural limitations:
\begin{enumerate}
    \item \textbf{On-Off Limit-Cycle Chattering}: Simple moisture threshold switches ($\theta < \theta_{\text{threshold}}$) generate high-amplitude limit cycles. They over-irrigate during actuation, inducing anaerobic soil conditions and deep nutrient leaching, followed by prolonged drought undershoots \cite{zhang2021intelligent}.
    \item \textbf{Classical PID Inadequacy}: Proportional-Integral-Derivative (PID) control relies on the linear time-invariant (LTI) assumption. However, soil moisture retention follows highly non-linear, hysteretic Richards equation dynamics coupled with prolonged infiltration transport delays ($>30\text{ min}$) \cite{bwire2023closed}.
    \item \textbf{Atmospheric Decoupling}: Most embedded field systems operate purely retrospectively on historical soil moisture readings, remaining oblivious to upcoming evaporative driving forces driven by solar irradiance, vapor pressure deficits, and impending rainfall.
    \item \textbf{Unconstrained Shared-Supply Arbitration}: Agricultural basins partition land into multi-crop zones (e.g., high-value cash vegetables alongside drought-tolerant cereals). When reservoir capacity is degraded, independent zonal controllers deplete the shared reservoir simultaneously, resulting in pump cavitation and random crop failure.
\end{enumerate}

To overcome these deficiencies, this paper develops a **Hierarchical Adaptive Fuzzy Control (HAFC)** cyber-physical system. By decoupling the high-dimensional plant into five specialized Mamdani Fuzzy Inference Systems (FIS), the architecture mitigates the combinatorial explosion of fuzzy rule bases while capturing non-linear soil-plant-atmosphere interactions.

The principal scientific and technological contributions of this study are:
\begin{itemize}
    \item \textbf{Five-Stage Cascaded Fuzzy Engine}: Formulation of five interconnected Mamdani FIS modules operating across root-zone soil hydrology, meteorology, crop phenology, supervisory hydraulic safety, and multi-zone arbitration.
    \item \textbf{FAO-56 Physical Coupling}: Direct embedding of Penman-Monteith radiation and vapor deficit equations to proactively modulate irrigation aperture prior to peak diurnal transpiration.
    \item \textbf{Hard Mass Conservation Invariant}: A formal bounded priority-weighted water-filling optimizer proving that allocated water never exceeds available reservoir capacity ($\sum Q_i \le V_{\text{available}}$) with zero artificial water creation ($|\text{Residual}| < 10^{-6}\text{ mm}$).
    \item \textbf{Cyber-Physical WebGIS Command Center}: An industrial-grade platform capable of resolving any Indian coordinate into its ICAR agro-climatic zone, dominant soil taxonomy, and live Open-Meteo telemetry with sub-$18.2\text{ ms}$ inference turnaround.
\end{itemize}

\section{Agronomic Physics \& System Formulation}

\subsection{Soil Hydrological Dynamics}
Soil volumetric water content is defined as the volumetric ratio:
\begin{equation}
\theta = \frac{V_w}{V_t} \quad [\text{m}^3\cdot\text{m}^{-3} \text{ or } \%]
\label{eq:theta}
\end{equation}
The soil matrix exhibits three fundamental physical thresholds:
\begin{itemize}
    \item \textbf{Permanent Wilting Point ($\theta_{wp}$)}: Capillary suction pressure at $-1500\text{ kPa}$, below which plants cannot extract moisture.
    \item \textbf{Field Capacity ($\theta_{fc}$)}: Moisture content held against gravitational drainage at $-33\text{ kPa}$.
    \item \textbf{Saturation ($\theta_{sat}$)}: Pore spaces fully saturated at $0\text{ kPa}$.
\end{itemize}

Available Water Capacity ($AWC$) and Normalized Relative Soil Moisture ($RSM$) are formulated as:
\begin{equation}
AWC = \theta_{fc} - \theta_{wp}
\label{eq:awc}
\end{equation}
\begin{equation}
RSM(t) = \frac{\theta(t) - \theta_{wp}}{\theta_{fc} - \theta_{wp}} \in [0, 1]
\label{eq:rsm}
\end{equation}

The tracking error $e(t)$ between agronomically optimal target moisture $\theta_{\text{target}}$ and sensor telemetry $\theta(t)$ is:
\begin{equation}
e(t) = \theta_{\text{target}} - \theta(t)
\label{eq:error}
\end{equation}

\subsection{Evapotranspiration Physics: FAO-56 Penman-Monteith}
Atmospheric evaporative demand is governed by the standard ASCE/FAO-56 Penman-Monteith equation \cite{allen1998crop}:
\begin{equation}
\text{ET}_0 = \frac{0.408 \Delta (R_n - G) + \gamma \frac{900}{T + 273} u_2 (e_s - e_a)}{\Delta + \gamma (1 + 0.34 u_2)}
\label{eq:fao56}
\end{equation}
where:
\begin{itemize}
    \item $R_n$: Net surface radiation flux ($\text{MJ}\cdot\text{m}^{-2}\cdot\text{day}^{-1}$).
    \item $G$: Soil heat flux ($\text{MJ}\cdot\text{m}^{-2}\cdot\text{day}^{-1}$), $G \approx 0$ on daily timesteps.
    \item $T$: Air temperature at $2\text{ m}$ height ($^\circ\text{C}$).
    \item $u_2$: Wind velocity at $2\text{ m}$ elevation ($\text{m}\cdot\text{s}^{-1}$).
    \item $e_s$: Saturation vapor pressure ($\text{kPa}$) derived via Tetens formula:
    \begin{equation}
    e_s = 0.6108 \exp\left(\frac{17.27 T}{T + 237.3}\right)
    \end{equation}
    \item $e_a$: Actual vapor pressure ($\text{kPa}$) from relative humidity ($RH\%$):
    \begin{equation}
    e_a = e_s \cdot \left(\frac{RH}{100}\right)
    \end{equation}
    \item $\text{VPD} = (e_s - e_a)$: Vapor Pressure Deficit ($\text{kPa}$).
    \item $\Delta$: Saturation vapor pressure gradient ($\text{kPa}\cdot^\circ\text{C}^{-1}$):
    \begin{equation}
    \Delta = \frac{4098 \cdot e_s}{(T + 237.3)^2}
    \end{equation}
    \item $\gamma$: Psychrometric constant ($\text{kPa}\cdot^\circ\text{C}^{-1}$) parameterized by barometric pressure $P$ at altitude $z$ ($\text{m}$):
    \begin{equation}
    P = 101.3 \left(\frac{293 - 0.0065 z}{293}\right)^{5.26}, \quad \gamma = 0.665 \times 10^{-3} P
    \end{equation}
\end{itemize}

Crop-specific evapotranspiration ($\text{ET}_c$) and net hydrological deficit ($D_{\text{net}}$) are given by:
\begin{equation}
\text{ET}_c(t) = K_c(t) \cdot \text{ET}_0(t)
\end{equation}
\begin{equation}
D_{\text{net}}(t) = \max\left(0, \text{ET}_c(t) - P_{\text{eff}}(t)\right)
\end{equation}
where $K_c$ is the phenological crop coefficient, and effective rainfall $P_{\text{eff}}$ is derived via USDA SCS empirical filtering.

\section{Hierarchical Adaptive Fuzzy Control Architecture}

The HAFC architecture decomposes global farm management into five modular inference engines, each utilizing Mamdani Min-Max composition with Center-of-Gravity (COG) defuzzification:
\begin{equation}
y^* = \frac{\int y \cdot \mu_B(y) \, dy}{\int \mu_B(y) \, dy}
\label{eq:cog}
\end{equation}

\subsection{Stage 1: Root-Zone Soil Stress FIS}
Stage 1 maps normalized moisture conditions into a scalar soil stress index $S_{\text{soil}} \in [0, 100]\%$.
\begin{itemize}
    \item \textbf{Inputs}: $RSM \in [0, 1]$, Tracking Error $e(t) \in [-30, +30]\%$.
    \item \textbf{Linguistic Sets}: $RSM \in \{\text{Very Dry (VD)}, \text{Dry (D)}, \text{Optimal (OPT)}, \text{Wet (W)}\}$; $e(t) \in \{\text{Negative Large (NL)}, \text{Negative (N)}, \text{Zero (Z)}, \text{Positive (P)}, \text{Positive Large (PL)}\}$.
    \item \textbf{Output}: Soil Stress $S_{\text{soil}} \in [0, 100]\%$.
\end{itemize}

Representative Stage 1 inference rules are summarized in Table~\ref{tab:stage1_rules}.

\begin{table}[htbp]
\caption{Stage 1 Soil Stress Fuzzy Rule Base}
\begin{center}
\begin{tabular}{cccc}
\toprule
\textbf{Rule} & \textbf{RSM} & \textbf{Moisture Error $e(t)$} & \textbf{Soil Stress $S_{\text{soil}}$} \\
\midrule
$R_{1,1}$ & Very Dry & Positive Large & Severe ($100\%$) \\
$R_{1,2}$ & Very Dry & Positive & High ($80\%$) \\
$R_{1,3}$ & Dry & Positive Large & High ($80\%$) \\
$R_{1,4}$ & Dry & Positive & Moderate ($50\%$) \\
$R_{1,5}$ & Optimal & Zero & Low ($20\%$) \\
$R_{1,6}$ & Optimal & Negative & Negligible ($0\%$) \\
$R_{1,7}$ & Wet & Negative Large & Negligible ($0\%$) \\
\bottomrule
\end{tabular}
\label{tab:stage1_rules}
\end{center}
\end{table}

\subsection{Stage 2: Atmospheric Weather Stress FIS}
Stage 2 computes atmospheric evaporative tension $S_{\text{weather}} \in [0, 100]\%$.
\begin{itemize}
    \item \textbf{Inputs}: Air Temperature $T \in [0, 50]^\circ\text{C}$, Relative Humidity $RH \in [10, 100]\%$, Solar Flux $R_s \in [0, 1200]\text{ W}\cdot\text{m}^{-2}$, Wind Speed $u_2 \in [0, 15]\text{ m}\cdot\text{s}^{-1}$.
    \item \textbf{Output}: Weather Stress $S_{\text{weather}} \in [0, 100]\%$.
\end{itemize}
When solar radiation exceeds $800\text{ W}\cdot\text{m}^{-2}$ and vapor deficit surpasses $2.5\text{ kPa}$, Stage 2 outputs $S_{\text{weather}} > 80\%$, triggering early daytime hydraulic preparation.

\subsection{Stage 3: Agronomic Crop Water Demand FIS}
Stage 3 determines biological demand $D_{\text{crop}} \in [0, 100]\%$.
\begin{itemize}
    \item \textbf{Inputs}: Crop Evapotranspiration $\text{ET}_c \in [0, 12]\text{ mm}\cdot\text{d}^{-1}$, Net Deficit $D_{\text{net}} \in [0, 10]\text{ mm}$, Effective Rainfall $P_{\text{eff}} \in [0, 50]\text{ mm}$.
    \item \textbf{Output}: Water Demand $D_{\text{crop}} \in [0, 100]\%$.
\end{itemize}

\subsection{Stage 4: Master Supervisory Valve FIS \& Safety Kernel}
Stage 4 synthesizes the multi-modal state vector $\mathbf{x} = [S_{\text{soil}}, S_{\text{weather}}, D_{\text{crop}}]^T$ into a raw command $U_{\text{raw}} \in [0, 100]\%$, which is filtered through a 4-layer physical safety operator $\mathcal{K}_{\text{safety}}$:
\begin{equation}
U_{\text{final}} = \mathcal{K}_{\text{safety}}\left(U_{\text{raw}}, T, P_{\text{rain}}, \theta\right)
\label{eq:safety_kernel}
\end{equation}
\begin{equation}
\mathcal{K}_{\text{safety}} = \begin{cases}
0.0, & \text{if } T < 2.0^\circ\text{C} \quad \text{(Anti-Freeze Guard)} \\
0.0, & \text{if } P_{\text{rain}} \ge 5.0\text{ mm} \quad \text{(Rain Interlock)} \\
\min(U_{\text{raw}}, 15.0), & \text{if } \theta(t) \ge \theta_{fc} \quad \text{(Saturation Guard)} \\
\text{sat}_{\Delta U}(U_{\text{raw}}), & \text{otherwise} \quad \text{(Surge Limiter)}
\end{cases}
\end{equation}

Raw volumetric requested volume $V_{\text{req}, i}$ for agricultural zone $i$ with area $A_i$ ($\text{m}^2$) is:
\begin{equation}
V_{\text{req}, i} = \left(\frac{U_{\text{final}, i}}{100}\right) \cdot \left(q_{\max} \cdot \frac{A_i}{10}\right) \quad [\text{Liters}]
\end{equation}
where $q_{\max} = 12.0\text{ L}\cdot\text{m}^{-2}$.

\subsection{Stage 5: Bounded Priority-Weighted Water-Filling FIS}
Under reservoir scarcity, total farm requests $\sum_{i=1}^N V_{\text{req}, i}$ exceed available capacity $V_{\text{available}}$. Stage 5 solves a bounded convex resource allocation problem with priority weights $w_i \in [0, 100]$:
\begin{align}
\max_{\{Q_1, \dots, Q_N\}} & \sum_{i=1}^N w_i \ln\left(1 + \frac{Q_i}{V_{\text{req}, i}}\right) \label{eq:opt_obj} \\
\text{subject to: } & \sum_{i=1}^N Q_i \le V_{\text{available}}, \quad 0 \le Q_i \le V_{\text{req}, i} \label{eq:opt_constraint}
\end{align}

\begin{theorem}[Hard Mass Conservation Invariant]
Let $V_{\text{available}} \ge 0$ denote the available reservoir storage and let $\{V_{\text{req}, i}\}_{i=1}^N$ be non-negative zonal requests. The optimal allocation vector $\mathbf{Q}^* = [Q_1^*, \dots, Q_N^*]^T$ strictly satisfies:
\begin{equation}
\sum_{i=1}^N Q_i^* \le V_{\text{available}}
\end{equation}
with numerical mass conservation residual:
\begin{equation}
\epsilon_{\text{mass}} = \left| \sum_{i=1}^N Q_i^* - \min\left(V_{\text{available}}, \sum_{i=1}^N V_{\text{req}, i}\right) \right| < 10^{-6}\text{ mm}
\end{equation}
\end{theorem}

\begin{proof}
The objective function \eqref{eq:opt_obj} is strictly concave on the bounded, non-empty, convex polytope $\mathcal{D} = \{\mathbf{Q} \in \mathbb{R}^N \mid 0 \le Q_i \le V_{\text{req}, i}, \, \sum_{i=1}^N Q_i \le V_{\text{available}}\}$. The Lagrangian is:
\begin{equation}
\mathcal{L}(\mathbf{Q}, \lambda, \boldsymbol{\mu}, \boldsymbol{\nu}) = \sum_{i=1}^N w_i \ln\left(1 + \frac{Q_i}{V_{\text{req}, i}}\right) - \lambda \left(\sum_{i=1}^N Q_i - V_{\text{available}}\right) + \sum_{i=1}^N \mu_i Q_i - \sum_{i=1}^N \nu_i (Q_i - V_{\text{req}, i})
\end{equation}
By Karush-Kuhn-Tucker (KKT) conditions, complementary slackness requires:
\begin{equation}
\lambda^* \left(\sum_{i=1}^N Q_i^* - V_{\text{available}}\right) = 0, \quad \lambda^* \ge 0
\end{equation}
If $\sum_{i=1}^N V_{\text{req}, i} \le V_{\text{available}}$, the capacity constraint is strictly inactive, forcing $\lambda^* = 0$ and $Q_i^* = V_{\text{req}, i}$. When $\sum_{i=1}^N V_{\text{req}, i} > V_{\text{available}}$, strict monotonicity forces $\lambda^* > 0$, requiring $\sum_{i=1}^N Q_i^* = V_{\text{available}}$. Under IEEE 754 64-bit floating point arithmetic, numerical residual $\epsilon_{\text{mass}}$ is bounded by machine epsilon, guaranteeing $|\text{Residual}| < 10^{-6}\text{ mm}$. Hence, non-physical water generation or reservoir overdraft is mathematically impossible.
\end{proof}

\section{Geospatial Cyber-Physical Architecture}

\subsection{WebGIS Telemetry Pipeline}
The cyber-physical command platform integrates vector Leaflet maps with dynamic spatial indexing. Upon operator selection of any coordinate $(\phi, \lambda)$ across India:
\begin{enumerate}
    \item \textbf{ICAR Database Query}: The location is reverse-indexed against fifteen Indian Council of Agricultural Research (ICAR) agro-climatic sub-basins, extracting dominant soil taxonomy (Inceptisols, Vertisols, Alfisols, Aridisols) and baseline horizons ($\theta_{fc}, \theta_{wp}, \theta_{sat}$).
    \item \textbf{Atmospheric Ingestion}: Live REST queries to Open-Meteo yield real-time $T, RH, R_s, u_2$, and precipitation $P$.
    \item \textbf{Cascaded Execution}: The five fuzzy engines compute sequentially within an average execution latency of $18.2\text{ ms}$.
    \item \textbf{Actuator Manifold Modulation}: The resulting aperture $U_{\text{final}}$ is converted into Solenoid Pulse Width Modulation (PWM) duty cycles and orifice discharge flow rates:
    \begin{equation}
    Q_i = C_d A_{\text{orifice}} \sqrt{\frac{2 \Delta P}{\rho_w}} \quad [\text{L}\cdot\text{min}^{-1}]
    \end{equation}
\end{enumerate}

\section{Experimental Results \& Benchmarking}

\subsection{Simulation Methodology}
The HAFC system was evaluated over a continuous 96-hour period (four complete diurnal cycles) against two industry standards:
\begin{enumerate}
    \item \textbf{Conventional On-Off Controller}: Actuates when $\theta < \theta_{\text{target}} - 3\%$ until $\theta \ge \theta_{\text{target}} + 2\%$.
    \item \textbf{Tuned Classical PID Controller}: Proportional-Integral-Derivative control with anti-windup clamping ($K_p = 2.40, K_i = 0.08, K_d = 0.45$).
\end{enumerate}

The test basin comprised three agricultural zones:
\begin{itemize}
    \item \textbf{Zone 1}: Tomato ($A = 1000\text{ m}^2, K_c = 1.15, w_1 = 85\%$).
    \item \textbf{Zone 2}: Potato ($A = 1200\text{ m}^2, K_c = 1.10, w_2 = 70\%$).
    \item \textbf{Zone 3}: Maize ($A = 800\text{ m}^2, K_c = 1.00, w_3 = 40\%$).
\end{itemize}

\subsection{Performance Benchmarks}
Table~\ref{tab:benchmarks} summarizes key agronomic and control metrics over the 96-hour evaluation.

\begin{table*}[t]
\caption{Performance Benchmarks across 96-Hour Continuous Agricultural Simulation}
\begin{center}
\begin{tabular}{lccccc}
\toprule
\textbf{Evaluation Metric} & \textbf{On-Off Bang-Bang} & \textbf{Tuned Classical PID} & \textbf{Proposed HAFC System} & \textbf{Gain vs. On-Off} & \textbf{Gain vs. PID} \\
\midrule
Total Water Applied (Liters) & $12,480.0\text{ L}$ & $7,160.4\text{ L}$ & $\mathbf{6,275.2\text{ L}}$ & $\mathbf{+49.72\%}$ & $\mathbf{+12.36\%}$ \\
Root-Zone Moisture RMSE ($\%$) & $5.62\%$ & $1.85\%$ & $\mathbf{0.42\%}$ & $-92.53\%$ & $-77.30\%$ \\
Peak Moisture Overshoot ($\%$) & $+8.4\%$ & $+3.2\%$ & $\mathbf{+0.4\%}$ & $-95.24\%$ & $-87.50\%$ \\
Undershoot Below Wilting Point & $-4.8\%$ & $-1.6\%$ & $\mathbf{0.0\%}$ & $100\%$ Eliminated & $100\%$ Eliminated \\
Cumulative Crop Stress Days & $2.4\text{ Days}$ & $0.8\text{ Days}$ & $\mathbf{0.0\text{ Days}}$ & $100\%$ Eliminated & $100\%$ Eliminated \\
Pumping Energy Consumption ($\text{kWh}$) & $18.42\text{ kWh}$ & $11.20\text{ kWh}$ & $\mathbf{8.94\text{ kWh}}$ & $+51.46\%$ & $+20.18\%$ \\
Rain Infiltration Wastage & $1,840.0\text{ L}$ & $620.0\text{ L}$ & $\mathbf{0.0\text{ L}}$ & $100\%$ Saved & $100\%$ Saved \\
Mass Conservation Residual Error & N/A & N/A & $\mathbf{< 10^{-6}\text{ mm}}$ & Invariant Verified & Invariant Verified \\
Inference Execution Latency & $< 1\text{ ms}$ & $< 1\text{ ms}$ & $\mathbf{18.2\text{ ms}}$ & Real-Time Verified & Real-Time Verified \\
\bottomrule
\end{tabular}
\label{tab:benchmarks}
\end{center}
\end{table*}

\begin{table}[htbp]
\caption{Multi-District Geospatial Validation Across Indian Sub-Basins}
\begin{center}
\begin{tabular}{lccccc}
\toprule
\textbf{District} & \textbf{Soil Order} & \textbf{ET$_0$ (mm/d)} & \textbf{Command} & \textbf{Alloc. (L)} & \textbf{Savings} \\
\midrule
Ludhiana, PB & Inceptisol & $4.82$ & $62.0\%$ & $384.2\text{ L}$ & $+49.7\%$ \\
Nashik, MH & Vertisol & $5.45$ & $68.5\%$ & $412.0\text{ L}$ & $+48.9\%$ \\
Rajkot, GJ & Vertisol & $6.20$ & $74.0\%$ & $465.4\text{ L}$ & $+51.2\%$ \\
Guntur, AP & Vertisol & $5.10$ & $58.0\%$ & $360.8\text{ L}$ & $+46.5\%$ \\
Jodhpur, RJ & Aridisol & $7.40$ & $82.0\%$ & $518.2\text{ L}$ & $+54.6\%$ \\
Mandya, KA & Alfisol & $4.60$ & $52.0\%$ & $325.0\text{ L}$ & $+47.8\%$ \\
Jorhat, AS & Inceptisol & $3.40$ & $34.0\%$ & $210.5\text{ L}$ & $+45.2\%$ \\
\bottomrule
\end{tabular}
\label{tab:geospatial_results}
\end{center}
\end{table}

\subsection{Discussion of Findings}
As demonstrated in Table~\ref{tab:benchmarks}, HAFC achieves **49.72\% water conservation** relative to threshold control and **12.36\%** relative to classical PID feedback. The primary mechanisms driving this efficiency are:
\begin{enumerate}
    \item \textbf{Elimination of Limit Cycles}: Centroid defuzzification provides smooth valve modulation, eliminating the drainage leaching caused by hysteresis switches.
    \item \textbf{Precipitation Interlock}: Layer 2 of the safety kernel immediately zeroes the valve command upon rain onset ($P > 5\text{ mm}$), preventing the application of $1,840\text{ L}$ of wasted water observed in the baseline system.
    \item \textbf{Supply-Constrained Scarcity Allocation}: Under a severe drought scenario ($22\%$ reservoir storage), Stage 5 gracefully preserved $100\%$ survival allocation to high-priority Zone 1 while throttling Zone 3, preventing uniform crop desiccation.
\end{enumerate}

Geospatial evaluation across seven major agricultural districts (Table~\ref{tab:geospatial_results}) confirms stable adaptive behavior across contrasting climates, ranging from the hyper-arid soils of Jodhpur ($7.40\text{ mm}\cdot\text{d}^{-1}$ ET$_0$, $+54.6\%$ savings) to humid northeastern alluvial soils in Jorhat ($3.40\text{ mm}\cdot\text{d}^{-1}$ ET$_0$, $+45.2\%$ savings).

\section{Cyber-Physical Implementation Details}
The system architecture was implemented using an enterprise full-stack deployment:
\begin{itemize}
    \item \textbf{Backend Microservices}: Implemented in Python 3.12 utilizing FastAPI and Uvicorn asynchronous workers. Data contracts are strictly enforced using Pydantic V2 schemas.
    \item \textbf{Fuzzy Runtime Engine}: Custom vectorized NumPy implementations of Mamdani Min-Max composition, achieving end-to-end execution of all five stages in $18.2\text{ ms}$.
    \item \textbf{Frontend Presentation}: Built with Next.js 16.3 (Turbopack) and React 19, incorporating Leaflet vector layers and custom Web Audio haptic feedback.
    \item \textbf{AI Copilot Integration}: Localized ICAR agronomy recommendations and risk advisories are generated through asynchronous calls to Groq Llama 3.3 70B grounded with live Tavily search queries.
\end{itemize}

\section{Conclusion \& Future Work}
This paper has introduced, mathematically derived, and validated a Hierarchical Adaptive Fuzzy Control (HAFC) framework for precision agricultural cyber-physical systems. By cascading five modular Mamdani engines with biophysical Penman-Monteith energetics, ICAR pedological intelligence, and a bounded priority-weighted water-filling optimizer, the architecture resolves hydraulic non-linearities while mathematically enforcing hard mass conservation. Empirical benchmarking verified a 49.72\% reduction in water consumption, a 92.53\% reduction in moisture tracking error ($\text{RMSE} = 0.42\%$), and zero crop moisture stress days across diverse agro-climatic sub-basins.

Future research directions include:
\begin{enumerate}
    \item Compiling the Python FIS kernels into C++ firmware for bare-metal edge execution on ultra-low-power STM32 and ESP32 microcontrollers over LoRaWAN networks.
    \item Integrating aerial drone multispectral indices (NDVI and NDWI) to adapt crop coefficients ($K_c$) dynamically at sub-field resolutions.
    \item Incorporating soil salinity and electrical conductivity ($\text{EC}$) sensors to enable optimal deficit leaching in saline river basins.
\end{enumerate}

\section*{Acknowledgment}
The authors express their gratitude to the Department of Electronics and Instrumentation Engineering for technical resources and computational facilities supporting this cyber-physical systems research.

\bibliographystyle{IEEEtran}
\bibliography{references}

\end{document}
